\section{Random merged pairs and their candidate graph}\label{sec:candidates}

The construction in Section~\ref{sec:deterministic-construction} calls for
two compatible split lines in every direction. We begin by supplying many
valid choices for each of these lines. We also record which choices
conflict: besides having large sets of candidates, we need their conflict
graph to have controlled degrees and, apart from a small specified set,
controlled codegrees.

Throughout this section, $M_g^0$ and $M_g^1$ are sampled independently and
uniformly from $\mathcal L_g$, independently for different directions
$g$. This sampling law allows the two constituents in a direction to
coincide. In that case they create no adjacency between distinct points
on a target line. With this convention the tuple validity tests of
Section~\ref{sec:deterministic-construction} apply throughout the sampling
space and agree with the original tests whenever every constituent pair
is distinct. We impose distinctness only after obtaining the estimates
below.

A \emph{successful tuple} in direction $g$ is a choice of four distinct
generators $d_0,d_1,d_2,d_3\ne g$ and bits
$\sigma_0,\ldots,\sigma_3$ such that
\[
 u_i=M_{d_i}^{\sigma_i}\cap M_{d_{i+1}}^{1-\sigma_{i+1}},
 \qquad i\in\mathbb Z/4\mathbb Z,
\]
are four distinct points on one $g$-line. Thus success concerns only
alignment and distinctness. A valid tuple must also pass the additional
tests in Section~\ref{sec:deterministic-construction}: its target line is
not a constituent, neither opposite pair is adjacent, and no external
point on the target line is adjacent to two designated positions.

For each line supporting valid tuples, retain one of them by a fixed
deterministic rule. Let $Q_g$ be the resulting set of tuples in direction
$g$. In particular, distinct members of $Q_g$ have distinct target lines.
Each of the two slots $(g,1),(g,2)$ carries a copy of $Q_g$. These copies
form the vertices of a graph: two vertices in different slots are joined
exactly when their tuples fail a compatibility condition from
Section~\ref{sec:deterministic-construction}; vertices in the same slot
are not joined. Choosing one vertex from each slot with no edges among
the chosen vertices is therefore precisely a compatible choice of the
required tuples. For a vertex $v$, write $L_v$ for its line, $F_v$ for its
four designated points, $\Gamma_v$ for its generator set, and $N(v)$ for
its graph neighborhood.

\begin{proposition}\label{prop:candidate-graph}
There are constants $c_*>0$ and $C_*<\infty$ such that, with probability
$1-o(1)$, the following assertions hold simultaneously, with
$\beta=1/8$ and $c_0=1/5$:
\begin{enumerate}[label=\textup{(\roman*)}]
\item $c_*q\le |Q_g|\le q$ for every direction $g$.
\item Every vertex has degree at most $C_*q$, and has at most $q^\beta$
neighbors in each slot.
\item For every $v$ there is a set $J(v)$ of at most $q^\beta$ vertices
such that $|N(v)\cap N(w)|\le q^\beta$ for $w\notin J(v)$.
Every member of $J(v)$ is adjacent to $v$ or is in its slot.
\item At most $\lceil\log q\rceil$ constituents pass through any affine
point, counting multiplicity. Every line other than one of the $M_g^1$
has at least $c_0q$ points outside $\bigcup_gM_g^1$.
\end{enumerate}
Consequently the event that these assertions hold and that
$M_g^0\ne M_g^1$ for every $g$ has probability $e^{-1}+o(1)$.
\end{proposition}

The proof has two counting stages. First we count successful tuples and
show that, after the validity tests, they occupy linearly many distinct
lines. We then bound incidences among tuples, their designated points,
and test lines; these estimates will give the graph bounds. A moment
argument supplies the uniform tail estimates at both stages. All counts
refer to the independent sampling law above, without conditioning on a
later selection of compatible split lines.

\subsection{A moment argument for small patterns}

When discussing exposure of the random choices, we call a direction a
\emph{colour}; its variable is the pair of random offsets specifying its
two constituents. A labelled pattern tests a bounded number of these
variables. Although different patterns may share colours, their counts
can be controlled if the same expectation bound holds after any bounded
exposure. The next lemma states the two forms we need. Constants may
depend on the fixed moment order and on the number of exposed colours,
but not on $q$ or on the values exposed.

\begin{lemma}\label{lem:pattern-moments}
Let $Z=\sum_a I_a$ count labelled patterns in independent colour variables,
where $I_a$ is the indicator of pattern $a$. For each label $a$, let $C_a$
be its deterministic set of tested colours, of uniformly bounded size,
so that $I_a$ is measurable in those variables. A bounded set $K_0$ of
colours may already have fixed values; all overlap tests below concern
the supports $C_a\setminus K_0$.
\begin{enumerate}[label=\textup{(\roman*)}]
\item If, after fixing any further bounded set of colours, the conditional
expected total is $O(1)$ uniformly, then $\E Z^s=O_s(1)$ for every fixed
positive integer $s$.
\item Suppose instead that the conditional expected total is $O(q)$,
and the conditional expected count of patterns using any of the further
exposed colours is $O(1)$. Overlap here excludes the initially fixed
colours. Then
\[
 \E|Z-\E Z|^{2s}=O_s(q^s).
\]
All expectations are conditional on the initially fixed colours, if any.
\end{enumerate}
\end{lemma}

\begin{proof}
For (i), expand $Z^s$ as a sum over ordered sequences of $s$ patterns.
After exposing all colours used by a prefix, the conditional expected
number of choices for the next pattern is $O_s(1)$, uniformly in the
values already seen. Summing one entry at a time proves the bound.

For (ii), work in the initially conditioned law and put $p_a=\E I_a$.
Expand the $2s$th centered moment as products of $I_a-p_a$. In each
sequence, form the graph whose entries are adjacent when their supports
outside $K_0$ share a colour. A term with an isolated entry vanishes,
because its centered factor is independent of all the other factors.
Every remaining sequence can be reordered so that each connected
component grows from its first entry, with every later entry sharing a
colour with an earlier one. Each component has at least two entries, so
at most $s$ entries start components.

To sum these terms in absolute value, bound each factor by $I_a+p_a$.
Given a prefix and its exposed values, an unrestricted next entry costs
$O_s(q)$, whereas an entry required to share a colour with the prefix
costs $O_s(1)$. The sums of $p_a$ obey the same bounds: for each fixed
set of earlier colour labels, average the uniform conditional bounds
over its possible values. Induction over the remaining entries therefore
gives $O_s(q^s)$ for each permitted ordering and component pattern.
There are only a bounded number of these, depending on $s$.
\end{proof}

We use part (i) for counts whose expectations remain bounded and part
(ii) for counts of order $q$ that must concentrate about their means.
Fix
\[
 \theta=1/100.
\]
By taking a sufficiently large fixed moment in (i), the probability that
such a count exceeds $q^\theta$ is smaller than any prescribed inverse
power of $q$. Part (ii) gives probability $O_s(q^{-s})$ for a deviation
by a fixed positive multiple of $q$. Both conclusions can consequently
be imposed over polynomially many tests.

We now verify the basic conditional counts for tuples. Fix a bounded
exposed colour set $K$. If a tuple has $j$ generators in $K$, let $f$
count its positions whose two incident generators are both exposed.
The cyclic arrangement of the four generators gives
\[
\begin{array}{c|ccccc}
j&0&1&2&3&4\\ \hline
f_{\max}&0&0&1&2&4.
\end{array}
\]
Every other position uses fresh offset variables, disjoint from those
used by the other nonfixed positions. Such a position is uniform in the
plane or on a fixed constituent whose direction differs from the target.
In either case its coordinate specifying a target parallel line is
uniform, independently of those of the other nonfixed positions.

For a fixed target direction, alignment therefore has probability at
most $q^{-3}$ when $f=0$, and at most $q^{-(4-f)}$ when $f\ge1$.
There are $O_K(q^{4-j})$ generator and bit assignments. Summing over
$j$ shows that the conditional expected number of successful tuples in
a fixed direction is $O_K(q)$, and that the count using an exposed
generator is $O_K(1)$. If the target line itself is prescribed, every
nonfixed position must lie on that line. We can then use
$q^{-(4-f)}$ also when $f=0$, giving conditional total $O_K(1)$ on
the prescribed line.

When all generators are exposed, we must also keep the distinctness
requirement in the definition of success. Four distinct fixed positions
can succeed in at most one target direction. We use this observation
whenever we sum fully exposed tuples over targets; otherwise such a sum
would introduce a spurious factor of $q$.

\subsection{Many distinct valid lines}

There are two obstacles to obtaining large sets $Q_g$: several successful
tuples may lie on the same line, and some successful tuples may fail
validity. We first control the concentration on individual lines, then
count the validity failures.

Let $H_g$ be the number of ordered successful tuples in direction $g$.
Before alignment is imposed, the four positions of a tuple are
independent uniform affine points: each uses two offsets of different
generators, and all eight offsets are used once. There are
$16(q)_4$ prescriptions, where $(q)_4=q(q-1)(q-2)(q-3)$, so
\[
 \E H_g=(16+o(1))q.
\]
Let $Y_g$ be the sum over $g$-lines of the square of the number of
successful tuples on that line. Equivalently, $Y_g$ counts ordered pairs
of successful tuples with a common target line. After bounded colour
exposure, choose the first tuple at expected cost $O(q)$; exposing its
colours then makes its target line specified, so the second tuple costs
$O(1)$. If the pair is required to use an already exposed colour, start
with a tuple having that overlap, at cost $O(1)$, and again count the
second on its specified line at cost $O(1)$.

These are exactly the conditional total and overlap bounds in
Lemma~\ref{lem:pattern-moments}(ii), for both $H_g$ and $Y_g$;
in particular, $\E Y_g=O(q)$. Its tail estimate and a union bound give,
with probability $1-o(1)$ simultaneously for every $g$,
\begin{equation}\label{eq:tuple-counts}
 H_g\ge8q,\qquad Y_g\le C_1q
\end{equation}
for an absolute constant $C_1$.

We next show that only $O(q^\theta)$ successful tuples in each direction
fail the extra validity tests, with probability $1-o(1)$ uniformly over
directions. For each failure, count the tuple together with a bounded
witness of that failure; overcounting is harmless. In each of the three
cases below the conditional expected witness count is $O(1)$ after any
bounded further colour exposure. Lemma~\ref{lem:pattern-moments}(i)
then supplies the required tail bound.

\paragraph{A constituent target line.}
First fix the two offsets of $g$. The target must be one of the two
specified constituent lines, and the prescribed-line estimate gives
conditional total $O(1)$.

\paragraph{An opposite adjacency.}
A witness consists of an extra colour and an orientation pairing two
opposite positions. This colour is neither the target nor a generator:
a generator pairs consecutive positions. If the witness colour is fresh,
first expose a successful tuple, at expected cost $O(q)$. The choices
for the witness then cost $O(q)q^{-2}=O(q^{-1})$, giving total $O(1)$.

If the witness colour is already exposed, distinguish whether a tuple
generator also overlaps the exposure. With such an overlap the basic
tuple count is $O(1)$. With all tuple generators fresh, enumerate the
target line instead. The exposed witness forces two exact positions on
that line, and the other two positions must lie on it. For fixed labels
and line this costs at most $q^{-6}$, against $O(q^5)$ choices of line
and tuple labels. This case also has conditional total $O(1)$.

\paragraph{An external point adjacent to two positions.}
Here the witness consists of two distinct extra colours, their
orientations, and their common external point $w$ on the target line.
Tuple generators cannot witness this condition. If both witness colours
are fresh, expose a successful tuple first. There are then $O(q^2)$
colour choices and $q$ choices for $w$, each with probability at most
$q^{-4}$. Their total cost is $O(q^{-1})$, against the $O(q)$ cost
of the tuple.

If a witness colour is exposed, first specify that colour and its
orientation. For a tuple overlapping the exposure, the tuple costs
$O(1)$, and $w$ is forced after the tuple's offsets are exposed. The
other witness costs $O(1)$ if it is exposed and $O(q^{-1})$ if it is
fresh. If instead all tuple generators are fresh, enumerate the target
line: the exposed witness prescribes $w$ and one exact tuple position.
For fixed labels and line the tuple costs at most $q^{-5}$, against
$O(q^5)$ choices. The remaining witness still costs at most $O(1)$.

These witnesses exhaust the extra validity failures. Subtracting their
counts from $H_g$ in~\eqref{eq:tuple-counts} leaves linearly many valid
tuples. Cauchy--Schwarz, together with $Y_g\le C_1q$, shows that these
tuples occupy at least $c_*q$ distinct lines. Retaining one tuple on
each such line gives $|Q_g|\ge c_*q$; the upper bound $|Q_g|\le q$
follows from the number of lines in a parallel class. This proves
Proposition~\ref{prop:candidate-graph}(i).

\subsection{Incidence estimates for candidates}

We have obtained enough candidates in every direction. To control their
conflicts, we now count tuples with prescribed point and line incidences.
All counts in this subsection include every successful tuple before
retention. They therefore bound the corresponding retained counts, and
their two slot copies change the bounds by at most a factor of two.

The following assertions hold simultaneously with probability $1-o(1)$:
\begin{enumerate}[label=\textup{(\alph*)}]
\item For a fixed target $g$ and any line $\ell$, at most $q^\theta$
successful tuples designate a point on $\ell$.
\item For any nonconstituent line $\ell$, at most $C_2q$ successful
tuples, over all targets, designate a point on $\ell$.
\item For any two distinct specified generators, at most $q^\theta$
successful tuples, over all targets, use both.
\item Each of the following counts, over all targets, is at most
$q^\theta$: tuples designating a specified point $y$; tuples whose line
passes through $y$ and which designate a point on a specified line $K$
not containing $y$; and tuples having two different designated positions
on two specified distinct nonconstituent lines, one on each.
\end{enumerate}
Assertions (a) and (b) control point-on-line conflicts within a slot and
over all slots, respectively. Assertion (c) controls conflicts from
shared generators. The more constrained incidences in (d) will control
common neighborhoods. To obtain simultaneous bounds, we verify the
conditional counts required by the moment lemma for each assertion.

For (a), a tuple with no exposed generator has probability at most
$4q^{-4}$: alignment costs $q^{-3}$, and a designated hit on the fixed
line costs a further $q^{-1}$. This remains true whether the hit forces
the target line or one point on it. Tuples overlapping the exposure have
conditional total $O(1)$ by the basic fixed-target bound. Thus the
bounded-total estimate of Lemma~\ref{lem:pattern-moments}(i) applies.

For (b), let $h$ be the direction of $\ell$ and first condition on the
offsets of $h$ avoiding $\ell$, as required for $\ell$ to be a
nonconstituent. Target direction $h$ itself is covered by (a). For the
other targets, let $j$ count generators in a bounded exposed set that
includes $h$. When $j=0$, the expected total is $O(q)$. When $j=1$,
there are $O(q^4)$ assignments including the target, and the hit
probability is at most $4q^{-4}$. To see the additional factor needed
here, condition on the common target-line coordinate. A fresh position
hits $\ell$ with probability $q^{-1}$, while a position on an exposed
constituent can hit $\ell$ for at most one common coordinate, since
that constituent is not $\ell$.

For $j=2,3$, the crude success probabilities $q^{-3},q^{-2}$ suffice
against $O(q^3),O(q^2)$ assignments. For $j=4$, the distinct fixed
positions allow at most one target. Hence all overlap counts are
$O(1)$. Tuples using $h$ satisfy the bounded-total estimate of
Lemma~\ref{lem:pattern-moments}(i). Those not using $h$ satisfy its
centered estimate under the initial conditioning on $h$, with mean
$O(q)$ uniformly. This proves (b), including a union bound over test
lines and the possible fixed offsets of $h$.

For (c), initially condition on the two specified generators. The
across-target estimates for $j\ge2$ in the preceding paragraph did not
use the requirement of a hit on $\ell$. They therefore give conditional
total $O(1)$ after any bounded further exposure, as required by
Lemma~\ref{lem:pattern-moments}(i).

For either of the first two tests in (d), choosing the target direction
prescribes the target line and one exact position on it. With $j=0$
the probability per assignment is at most $4q^{-5}$, and with $j=1$
it is $O(q^{-4})$. In the latter case, even when the exact position is
constrained to an exposed constituent, specifying it costs at least
$q^{-1}$; each of the other three positions must lie on the prescribed
target line. The across-target bounds for $j\ge2$ suffice as before.

For the last test in (d), condition on the offsets in the direction or
directions of the two test lines, avoiding those lines. For a target
different from these directions, requiring two different positions to
meet the two lines costs two further factors $q^{-1}$ when $j=0$, and
at least one further factor when $j=1$. Use the crude success bounds
when $j\ge2$. If the target equals a test direction, its line must be
the specified test line; the prescribed-line estimate then suffices.
Thus every test in (d) has conditional total $O(1)$.

There are only polynomially many tests in (a)--(d). Arbitrarily high
fixed moments and a union bound prove the assertions. All estimates made
after initial conditioning are uniform in the allowed fixed offsets,
so averaging over those offsets preserves the bounds.

\subsection{Degrees, codegrees, and initial occupancy}

We now translate the incidence estimates into the graph assertions of
Proposition~\ref{prop:candidate-graph}. For the first compatibility
condition, the two ways a conflict can occur are represented by
\[
 B_v=\{w:L_w\cap F_v\ne\varnothing\},\qquad
 E_v=\{w:F_w\cap L_v\ne\varnothing\}.
\]
The one-per-line retention rule gives $|B_v|=O(q)$ and at most four
members of $B_v$ in any slot. Since a valid tuple's line $L_v$ is not
a constituent, (b) gives $|E_v|=O(q)$; (a) gives at most $q^\theta$
members of $E_v$ in any slot.

The other compatibility conditions contribute fewer conflicts. Two
shared generators give $O(q^\theta)$ vertices by (c). For the condition
on slots with the same target direction, every merged pair witnessing
a conflict meets $F_v$. Once the constituent multiplicity estimate
below holds, there are at most $4\lceil\log q\rceil$ such pairs.
Assertion (a), applied to their constituents in the other slot of this
direction, bounds the candidates they meet. These conflicts contribute
$O(q^\theta\lceil\log q\rceil)$ vertices. Together the three types
give the required degree bounds.

For codegrees, the exceptional set has a direct geometric description.
Define $J(v)$ to consist of vertices on the same line as $v$ or sharing
a designated point with $v$. The retention rule and the first test in
(d) give $|J(v)|=O(1+q^\theta)$. Every member is either in $v$'s
slot or adjacent to $v$, as asserted in part (iii).

Fix $w\notin J(v)$. Then $L_v,L_w$ are distinct and $F_v,F_w$ are
disjoint. We first bound common neighbors arising from the first
compatibility condition, by considering the four intersections of their
$B$ and $E$ sets. A member of $B_v\cap B_w$ has a line joining a
point of $F_v$ to a point of $F_w$. There are only a bounded number of
such lines, hence a bounded number of candidate vertices by retention.

For $B_v\cap E_w$, choose $y\in F_v$ on the candidate line. If
$y\notin L_w$, the point-line test in (d) applies. If $y\in L_w$,
the candidate line is either $L_w$, or its designated point on $L_w$
must be $y$. The one-per-line rule and the point test in (d) handle
these two possibilities. The same argument bounds $E_v\cap B_w$.
Finally, a member of $E_v\cap E_w$ either has two different designated
positions on $L_v,L_w$, or designates their intersection. The last or
first test in (d), respectively, applies.

The remaining conflict types add at most
$O(q^\theta\lceil\log q\rceil)$ to these common-neighbor bounds,
by the bounds already used for degrees. Since $\theta<\beta=1/8$,
the exceptional-set size, codegree bounds, and within-slot degree bounds
are all at most $q^\beta$ for sufficiently large $q$. This proves
parts (ii) and (iii), subject only to the multiplicity estimate we now
establish along with the rest of part (iv).

The number of constituents through a fixed affine point has distribution
$\operatorname{Bin}(2(q+1),1/q)$, counting multiplicity. Its tail at
$\lceil\log q\rceil$ is $q^{-\omega(1)}$, so a union bound over
the $q^2$ affine points proves the multiplicity estimate.

For the uncovered-point estimate, fix a line $\ell$ of direction $h$.
The $q$ lines $M_g^1$ with $g\ne h$ hit $\ell$ at independent
uniform points. If $U_\ell$ counts points missed by these hits, then
\[
 \E U_\ell=q(1-1/q)^q.
\]
Changing one hit changes $U_\ell$ by at most one. Apply
Lemma~\ref{lem:independent-trial-jumps} to these $q$ independent trials
with $L=1$, and then Lemma~\ref{lem:bounded-drift-concentration} to
their exposure martingale. Since $(1-1/q)^q\to e^{-1}>1/5$, there is
a fixed $\delta>0$ such that $\E U_\ell-q/5\ge\delta q$ for all
sufficiently large $q$. Taking deviation $u=\delta q$, a fixed positive
multiple of $q$, gives probability $\exp(-\Omega(q))$ that
$U_\ell<q/5$. If $\ell\ne M_h^1$, the remaining parallel line
$M_h^1$ does not meet it. A union bound over all lines therefore proves
the uncovered-point assertion in (iv).

All failure probabilities above sum to $o(1)$ in the independent
sampling model. Independently across directions, the probability that
each constituent pair is distinct is
\[
 (1-1/q)^{q+1}=e^{-1}+o(1).
\]
Intersecting this event with the event of probability $1-o(1)$ just
proved gives the final assertion of
Proposition~\ref{prop:candidate-graph}.

